\section{Related work}
\label{sec:related}

\subsection{Fish and Literature specifically}

The literature on this game is human and informal. McLeod's Pagat entry
\cite{pagat} is canonical for the rules --- including our 54-card nine-set deck and
the variant in which any declaration error awards the half-suit to the opponents
--- and for a tactics list covering deliberate misses that inform a partner,
locking out a dangerous opponent by never asking them anything, and exhausting one
rank across all opponents before switching rank; it also records Ali Salahuddin's
partner-signalling convention and Guy Srinivasan's Forced Claims rule. Develin's
chapter \cite{develin} is the deepest strategy writing we found and describes
exactly our dialect: the blackball example, the rule against declaring a half-suit
one holds entirely, the observation that asking teaches opponents as much as
teammates --- so asks inside a team-owned half-suit are a free channel --- and the
corpus's only quantitative endgame, where the alternative to a coin-flip
declaration would need success probability $3/2$. He forbids pre-agreed
conventions outright, contradicting Pagat. Wikipedia
\cite{wikiliterature} names the ``stalemate breaker''; secondary pages
\cite{dorsa,strategy-site} restate blackballing, asking the asker, and an endgame
delegation heuristic.

Computational prior art is close to absent. The review corpus reports that arXiv
searches on 21 August 2026 for \texttt{"Canadian Fish"},
\texttt{"Literature card game"} and \texttt{"half-suit"} each returned nothing,
that a GitHub sweep found no repository containing search, planning or a working
learner, and that Literature is not among RLCard's environments \cite{rlcard}. The
one real agent is Somani's \texttt{literature} \cite{somani,somani-server}: about
$1{,}940$ lines of Python with a complete rules engine and a sound deduction fixed
point. It plays the 48-card, eight-half-suit ``remove the sevens'' variant and
trains only on four players, so the six-player game's hardest sub-problem is
missing --- with one teammate a half-suit allocation spans $2^6$ possibilities and
is usually forced, with two it spans $3^6 = 729$. Its learner, a
$1149 \to 100 \to 1$ perceptron fitted for $1{,}030{,}945$ iterations, never
bootstraps: it regresses a hand-written score of $\pm 20$ per ask outcome plus
$\pm 100$ at the end, which punishes the deliberate miss both Pagat and Develin
recommend, and its move generator forbids known misses except as a fallback. A
verified bug compounds this --- team identity is an integer and the winner a plain
enumeration member, so the terminal comparison is always false and every move in
every game received $-100$. The published model never saw a win/loss gradient and
is functionally a greedy maximum-likelihood asker. Declarations fire the instant a
half-suit is provably owned, wrong claims are unpenalised by construction, and no
strength numbers have ever been published. The other public repositories are user
interfaces; one closed-source Android application advertises four- and six-player
bots with no stated methodology (unverified) \cite{playstore}.

\subsection{Search under imperfect information}

The standard recipe for hidden-hand card games is determinization, or Perfect
Information Monte Carlo: sample an assignment of the hidden cards consistent with
the public history, solve the resulting perfect-information game with a
double-dummy solver, and play the best average move. Levy proposed it for Bridge
(a source the corpus could not verify) \cite{levy}; Ginsberg's GIB made it work,
reweighting sampled deals by Bayes and finding fifty worlds sufficient
\cite{ginsberg-gib}. Frank and Basin showed the method unsound however many worlds
are drawn, naming \emph{strategy fusion} --- the searcher plays a different
strategy in each world, although a real player must commit to one per information
set --- and \emph{non-locality}, where a node's value depends on parts of the tree
outside its subtree because informed opponents steer play away from it
\cite{frank-basin-ai}; the repairs that work are those attacking non-locality too
\cite{frank-basin-aaai}. Long, Sturtevant, Buro and Furtak explained when
determinization succeeds at all, through leaf correlation, bias and a
disambiguation factor, placing Skat and Hearts where PIMC is near-optimal and Kuhn
poker where it is not \cite{long-pimc}. Ginsberg names the structural cost:
because each determinization assumes the missing information will arrive at the
next decision, PIMC never makes an information-gathering play.

Information Set MCTS instead builds one tree over information sets, redeterminizes
each iteration, and explores on the count of iterations in which a move was
\emph{available} rather than on parent visits \cite{cowling-ismcts}. Its variants
coincide when all moves are fully observable, as in Fish, and one tree of ten
thousand iterations beat forty trees of two hundred and fifty by $22.9\%$;
re-determinization at non-root seats and self-determinization extend it
\cite{goodman-ris,cowling-bluffing}, though ISMCTS is unsound too, its
exploitability sometimes \emph{increasing} with computation \cite{lisy-oos}. The
Skat and Bridge engineering lines matter as much: Kermit's inference module alone
was worth $+217$ tournament points per 36 games \cite{buro-kermit}, and a Spades
post-mortem traced blunders to a sampler placing a card of true probability $1/3$
into $87\%$ of its worlds \cite{whitehouse-spades}.

None of this is the right foundation here. In Fish the \emph{perfect-information}
game is degenerate: a clairvoyant player only ever names cards the target actually
holds, so never fails an ask and never loses the turn, and declares each half-suit
correctly the instant their team completes it. The double-dummy value of almost
every position is ``the team on turn takes everything it can reach''; the solver
barely discriminates among root moves, and averaging over sampled worlds leaves
\[
  \arg\max_{(t,c)} \; \mathbb{E}\bigl[V^{\mathrm{PI}}\bigr]
  \;\approx\; \arg\max_{(t,c)} \; \Pr[\,t \text{ holds } c \mid h\,],
\]
a one-ply heuristic blind to the information an ask leaks, the information a
refusal supplies, the option value of postponing a declaration, and partner
signalling. A Fish double-dummy solver will be easy to write, fast, and worthless
as an engine --- useful only for hit-probability features and as a sparring
baseline. The corpus's reading of \cite{long-pimc} sharpens the point: Fish should
have a disambiguation factor far above the $\approx 0.6$ measured for Skat and
Hearts --- nominally PIMC's best regime --- yet its pre-terminal leaf structure
sits in the corner where almost every move wins in the sampled world, which that
paper never measured and dismissed as ``very boring games''. Fish is not boring,
precisely because of the imperfect information determinization discards. That is
why we build an exact belief engine (\S\ref{sec:belief}) and a decision-theoretic
declaration rule (\S\ref{sec:theory}) instead. One smaller finding survives: Go
Fish sits at the maximum of Zuckerman, Felner and Kraus's Opponent Impact measure
\cite{zuckerman-mpmix}, so which opponent to ask is a first-class argument here,
not a tie-break.

\subsection{Equilibrium computation and team games}

The equilibrium line runs from counterfactual regret minimisation
\cite{zinkevich-cfr} through sampling variants \cite{lanctot-mccfr}, CFR+
\cite{tammelin-cfrplus} and discounted CFR \cite{brown-dcfr} to function
approximation \cite{brown-deepcfr,steinberger-dream}. CFR+ does relatively badly
where some actions are very costly mistakes \cite{brown-dcfr}, and a wrong
declaration is exactly that; DREAM's exploration term is exponential in depth, and
a Fish game runs sixty to a hundred and twenty decisions. ReBeL recasts an
imperfect-information game as a continuous-state perfect-information game over
\emph{public belief states} \cite{brown-rebel}, building on the common-information
reduction of decentralised control \cite{nayyar-common} and the public-belief
construction used in Hanabi \cite{foerster-bad}; Student of Games extends it and
names Scotland Yard as the closest published public-observation domain
\cite{schmid-sog}. Fish is an \emph{exact} public-belief-state game, which is what
\S\ref{sec:belief} exploits.

Three teammates share a payoff but hold different information and cannot
communicate, so the team is not one player with pooled knowledge --- assuming
otherwise is a second flavour of strategy fusion --- and Nash equilibrium over
individual strategies is not the target either. The right concept is the
team-maxmin equilibrium with correlation (TMECor): agree on a correlated joint
strategy before the deal, then execute it without communicating. Celli and Gatti
give the complexity --- TMECor is FNP-hard, the team best-response oracle APX-hard,
and the worst-case price of uncorrelation grows with the leaf count $|L|$, with
reported bounds of $|L|/2$, $|L|/4$ and $\sqrt{|L|}$ \cite{celli-gatti}. Farina
and coauthors connect ex-ante team collusion to extensive-form correlation
\cite{farina-collusion}, and the team belief DAG and two-player conversion make
the team a single coordinator issuing prescriptions, giving a two-player zero-sum
game whose equilibria are realization-equivalent to TMECor, at the cost of a
representation exponential in how much private information distinguishes teammates
within a public state \cite{zhang-tbdag,carminati-tpi}. Learning-based attacks
exist \cite{mcaleer-teampsro,xu-fxp}; the latter proves self-play converges to
\emph{local} team equilibria --- the formal version of the complaint that a bot's
asks stop meaning anything.

We attempt none of this: one seat's information set at deal time holds
$45!/(9!)^5 \approx 1.9 \times 10^{28}$ deals, and a prescription would have to
name an action for each teammate at each of their possible hands. The cost should
be stated plainly. Nothing here is an equilibrium result --- win rates are measured
against a fixed population, robustness is a worst-case-over-panel claim rather
than an exploitability bound, and the fitted exploiter of \S\ref{sec:protocol} is a
lower-bound proxy whose failure mode is known: a low score against an exploiter is
not evidence of low exploitability \cite{lisy-lbr}.

\subsection{Cooperative play through public actions}

Hanabi is the reference benchmark for coordination through observable actions
\cite{bard-hanabi}, and is relevant because Fish likewise has no side channel. The
Bayesian Action Decoder made the public belief explicit, reached $24.174$ of $25$,
and attributed roughly $40\%$ of its information transfer to learned conventions
rather than grounded content \cite{foerster-bad}. The Simplified Action Decoder
showed that letting teammates condition on the \emph{greedy} rather than the
exploratory action is a prerequisite for conventions to survive
$\varepsilon$-greedy training \cite{hu-sad}. Other-Play exposed how much of a
self-play score is convention rather than skill: the same agents scored $23.97$
with themselves and $2.52$ paired across independent runs \cite{hu-otherplay}.
Off-Belief Learning is an explicit dial on convention depth, from a grounded
level-one policy near $20.9$ to $24.10$ at level four, with a uniqueness result
that makes independent runs interoperable \cite{hu-obl}; test-time search alone
lifted a blueprint from $24.08$ to $24.61$ \cite{lerer-sparta}. Bridge bidding is
the other model system, an auction that is literally a learned language: bidding
networks beat a strong commercial program by $0.25$ IMP per board
\cite{rong-bidding}, joint policy search --- which scores simultaneous changes at
several information sets, exactly what inventing a convention requires --- improved
a system from $+0.29$ to $+0.63$ IMPs per board \cite{tian-jps}, and an agent
trained with human partners gained $+0.97$ IMPs per deal \cite{lockhart-bidding}.

The asymmetry separating Fish from both is the audience. Hanabi has no adversary,
so every bit of signal is pure profit --- which is why its conventions are so
elaborate --- and in bridge the harm a leaked auction does is diffuse. In Fish a
leaked card location is directly executable: once it is public that a player holds
a card, any opponent holding another card of that half-suit takes it with
probability one and keeps the turn. Conventions here carry a genuine eavesdropping
cost. The corpus states the sharpest version information-theoretically: because
teammate and opponent hands are exchangeable given the public history, an absolute
codebook --- one whose meaning is fixed by public information alone --- is decoded
identically by all five listeners, so with two teammates against three opponents
the Wyner-style secrecy rate of such a convention is non-positive
\cite{wyner,csiszar-korner}; only a receiver-relative codebook, whose semantics
depend on the receiver's own holdings, has positive rate, at the price of
ambiguity at the encoder. We report this as an argument, not a theorem: the corpus
gives a structural claim with an informal exchangeability justification and no
proof, and we have not proved it either. It is the design rationale for building
v0.4 with no private channel at all. Public-signalling economics predicts the same
outcome, cheap talk separating before two audiences only when credible against
both and otherwise collapsing toward pooling --- the regime Farrell and Gibbons
call subversion \cite{farrell-gibbons}. The Hanabi construction that most clearly
fails to port is the modular hat-guessing code, which works because each receiver
sees every other receiver's hand and can subtract it \cite{cox-hanabi}; in Fish
nobody sees any hand. Any claim that an agent has learned a convention must also
survive the scramble ablation of \cite{lowe-pitfalls}, where high
speaker--listener consistency coexisted with zero causal effect.

\subsection{Inference over deals}

Counting the deals consistent with a public history is a permanent computation.
Ryser's inclusion--exclusion and Glynn's sign-vector formula are exact at
$O(2^n n)$ \cite{ryser,glynn}; at $n = 54$ that is about $1.8 \times 10^{16}$, so
they serve only as unit tests on tiny instances. The Jerrum--Sinclair--Vigoda
chain is a fully polynomial randomised approximation scheme \cite{jsv}, but a
practical study measured it at $50{,}634$ seconds for $n = 10$ against Ryser's
$0.003$ seconds \cite{newman-vardi}. Matrix scaling is the usual approximation ---
Sinkhorn iteration, with the deterministic strongly polynomial analysis bounding
the permanent within a factor $e^n$ \cite{lsw} --- and it is what our fast path
uses (\S\ref{sec:fast}); the corpus measures its cost as two to six percent
maximum marginal error on a Fish-sized instance, worst exactly where a real
position gets hard, reaching $5.9\%$ at $60\%$ void. The contingency-table
literature contributes sequential importance sampling \cite{chen-sis}, the warning
that it can underestimate by an exponential factor under any row or column
ordering while appearing to have converged \cite{bezakova-sis}, and, closest in
spirit to what Fish permits, polynomial approximate counting of tables with a
constant number of rows \cite{cryan-dyer}. Constrained deal generation in practice
is done by rejection or by learning: GIB deals against suit-length restrictions and
then discards deals whose bids its own bidding module would not have made ---
rejection sampling against a policy, one to two seconds for fifty worlds
\cite{ginsberg-gib,pavlicek} --- while Skat engines enumerate and reweight with a
learned per-card location model or policy-based reach probabilities
\cite{solinas-supervised,rebstock-inference}. The general problem is hard:
constructing a history consistent with a public state is FNP-complete in the worst
case, and enumeration is polynomial only for \emph{sparse} public trees
\cite{solinas-history}.

Fish is unusually favourable. Cards move only by publicly observed transfer, so
the hidden state never grows --- it is the initial deal and nothing else --- and
each player's share of the initial deal stays constant at nine all game. The counting matrix
therefore has at most six distinct column types, one per seat, and its permanent
becomes polynomial rather than \#P-hard: a dynamic program over the capacity
vector costs about $5 \times 10^5$ multiply--adds from a player's seat. Two
caveats deserve recording: tractability needs the constant column count \emph{and}
capacities that are small numbers written in unary, since counting contingency
tables is \#P-complete even with two rows once margins are binary-encoded; and ask
legality is disjunctive, which would normally destroy the product form --- the
escape being that every legality certificate is confined to a single half-suit, so
exhaustive enumeration within a block keeps it exact. A reference implementation
with nine blocks and 27 exact legality constraints runs in $0.72$ seconds of
unoptimised NumPy, milliseconds once ported to C. Nothing comparable holds in
Bridge, where the constraints are soft auction inferences and no such collapse of
column types occurs.

\subsection{Evaluation methodology}

Card-game results are dominated by deal variance, so the standard remedy is a
matched design: duplicate bridge replays each board with the same cards in the
same seats, and the computer poker competition adds common seeds across pairings,
reportedly cutting the standard deviation to about two thirds --- a figure the
corpus could verify only partially \cite{acpc}. The resampling unit is then the
deal, not the game, since replays of one deal form one correlated cluster; and the
raw percentile bootstrap should be avoided, having been measured to fail at $5.7$
to $11.2\%$ against a nominal $5\%$ \cite{jordan-rleval}. Poker also supplies
control variates: advantage-sum estimators \cite{zinkevich-divat}, importance
sampling over imaginary observations \cite{bowling-is}, MIVAT's least-squares value
function \cite{white-mivat}, and AIVAT, which cut standard deviation by $68.8\%$ in
heads-up no-limit hold'em and $99.9\%$ in Leduc with full knowledge
\cite{burch-aivat,avaivat}. The gains are domain-sensitive: MIVAT delivered only
about eighteen to twenty percent in six-player poker, the closest analogue to our
setting.

For absolute quality the tools are best-response computation and its relaxation,
local best response, which showed every entrant in the 2016 computer poker
competition to be exploitable for more than $3{,}180$ milli-big-blinds per hand
\cite{lisy-lbr}. Two lessons carry over verbatim: it gives only a lower bound, so a
low score proves nothing --- it lost $536$ against a bot whose true exploitability
in the same action set was $90$ --- and the point at which the exploiter may act
must be swept, since restricting it to later rounds raised measured exploitability
by nearly an order of magnitude. Refining an abstraction likewise does not
monotonically reduce exploitability \cite{waugh-abstraction}. For ranking a
population, Bradley--Terry-style models \cite{coulom-whr,herbrich-trueskill} fail
twice over when that population is a training run's own checkpoints: they are not
redundancy-invariant, so duplicating one agent redistributes everyone's rating,
and they cannot represent cycles. Balduzzi and coauthors offer a maximum-entropy
Nash average over the antisymmetric logit matrix instead
\cite{balduzzi-reeval,omidshafiei-alpharank}.

Three documented traps shaped our protocol. Monitoring an interval computed for a
fixed sample size and stopping when it looks good is not a test: in simulation it
produced a $61.35\%$ false-positive rate under the null \cite{avaivat}, which is
why sequential designs on normalised Elo exist \cite{nelo}. Resampling games rather
than deals inflates apparent precision whenever several games share a deal. And
fitting a variance-reduction heuristic on the evaluation data destroys the
guarantee that motivated it: AIVAT is unbiased for \emph{any} value function,
which constrains its expectation and not the realised estimate, and
\cite{kim-sandholm} tuned the heuristic on a published ten-thousand-hand
six-player dataset until all fourteen players simultaneously appeared to be large
winners --- which the zero-sum structure makes impossible. Accordingly
\S\ref{sec:protocol} uses six-rotation duplicate blocks, resamples deals and not
games, freezes its held-out seed banks before the configuration is fixed, and uses
a fitted exploiter as its only exploitability proxy. We deliberately do not use
AIVAT or MIVAT: a learned control variate fitted by us, on our own agent, is the
exact pathology this literature warns about.
